% Resumo (NBR 6028). Números iguais aos dos capítulos (situação de 17/11/2026).
% "56 histórias": número de linhas da Tabela tab:backlog (gerado/backlog.tex),
% igual à soma da Tabela tab:moscow-resumo (11 + 13 + 24 + 8).
% Palavras-chave: "lavagem de dinheiro" entrou no lugar de "contratações
% públicas", para seguir o título (decisão do autor, 17/11/2026).
% PENDENTE: atualizar a frase sobre as tipologias de lavagem quando a
% implementação em andamento terminar.
% PENDENTE: atualizar os números se R2 mudar alguma regra das revisões, e a
% frase final se a avaliação com usuários for executada.
\begin{resumo}
 Este trabalho desenvolve um protótipo de buscador de inteligência de fontes
 abertas (OSINT) para a polícia judiciária investigar corrupção e lavagem de
 dinheiro com dados abertos brasileiros, a partir de contratações públicas,
 empresas, candidaturas e bens declarados. Os requisitos partiram de um levantamento com
 especialista e foram confrontados com a legislação e com uma revisão
 sistemática da literatura (RSL 1). Uma segunda revisão (RSL 2) orientou o
 processo de desenvolvimento. As duas revisões são relatadas segundo o PRISMA
 2020 e o SEGRESS. A RSL 1 trata de ferramentas OSINT para órgãos policiais e
 de inteligência. Dos
 1.928 registros recuperados, restaram 1.139 após a remoção de duplicatas, e
 57 estudos primários foram incluídos, 23 da busca nas bases e 34 do
 \textit{snowballing}. Nenhum estudo é do Brasil ou de outro país da América
 Latina, nenhum trata de fraude ou corrupção em contratações públicas,
 nenhum tem a lavagem de dinheiro como objeto, e só 7 tratam explicitamente
 da cadeia de custódia. A RSL 2 trata dos
 efeitos de assistentes e agentes baseados em modelos de linguagem de grande
 porte no desenvolvimento de software. Dos 8.181 registros, restaram 4.323
 após a remoção de duplicatas, e 130 relatos, correspondentes a 129 estudos,
 foram incluídos. Em produtividade, 51 dos 89 estudos que mediram o efeito
 relatam direção positiva, e a direção é mais favorável quando o efeito é só
 percebido. Em segurança, 14 dos 23 estudos que mediram o efeito relatam
 direção negativa. A comparação com sete ferramentas existentes, feita pela
 documentação oficial, não encontrou nenhuma que reúna fontes abertas
 brasileiras de contratações, tipologias de fraude e preservação da prova
 para o inquérito. Os requisitos foram escritos como 56 histórias de usuário,
 priorizadas pelo método MoSCoW. O protótipo, de código aberto, tem coleta só
 passiva, cadeia de custódia com \textit{hash} SHA-256 e carimbo de tempo em
 cada captura, trilha de auditoria só de acréscimo, dez tipologias de fraude
 em contratações por regras explícitas, cujos alertas são pistas que o
 investigador confirma ou descarta, e extração de entidades por IA opcional
 e local. As tipologias de lavagem com dados abertos, que só dão indícios
 cadastrais, e a importação separada de dados bancários sob ordem judicial,
 só com dados sintéticos, estão no \textit{backlog}. Foi
 desenvolvido com agentes de IA, aplicando parte das práticas candidatas da
 RSL 2, e tem
 342 testes automatizados no \textit{back-end} e 234 no \textit{front-end}.
% PENDENTE: atualizar as contagens de testes ao fim das correções em andamento.
 Os resultados das revisões são provisórios, porque etapas feitas pela IA
 como terceira opinião ainda aguardam a confirmação dos dois revisores. O
 protótipo tem testes automatizados com dados sintéticos e teve os conectores
 conferidos contra as fontes reais, mas ainda não foi avaliado com policiais. Conclui-se que a literatura
 e as ferramentas analisadas não cobrem a combinação de domínio de
 contratações públicas e preservação da prova, que a literatura revisada não
 trata de lavagem de dinheiro, e que o protótipo mostra que essa combinação
 pode ser implementada.

 \vspace{\onelineskip}

 \noindent
 \textbf{Palavras-chave}: inteligência de fontes abertas. investigação
 criminal. lavagem de dinheiro. cadeia de custódia. revisão sistemática.
\end{resumo}
